
\begin{proposition}[Credited re-derivation]\label{cdp:st:D.1.3}

\textbf{Exact check on the retained construction.}
In the ordered homology basis
$(\widehat\gamma,\widehat u,\widehat w,\widehat\delta)$,
the two finite monodromies are
\[
A_1=\begin{pmatrix}1&0&0&0\\6&0&1&0\\-6&-1&-1&0\\-2&1&0&1\end{pmatrix},
\qquad
A_2=\begin{pmatrix}1&0&0&0\\0&0&-1&0\\-6&1&0&0\\3&0&1&1\end{pmatrix}.
\]
They satisfy $A_1^3=A_2^4=1$; the third generator is
$(A_1A_2)^{-1}$. The columns of
$[A_1-I\mid A_2-I]$ span the primitive subgroup
$\ker(\gamma:\Lambda\to\Z)$, where $\gamma=(1,0,0,0)$.
Indeed its first row is zero and its Smith diagonal is $(1,1,1,0)$.
The invariant homology vectors are $\Z\widehat\delta$, while the
invariant cohomology covectors are $\Z\gamma$. Consequently
\[
\Lambda_{\langle A_1,A_2\rangle}=\Z,
\qquad t'=1,
\qquad r=3-1+1=3.
\]
The printed count would give $r=3-1+4=6$.
These assertions are proved by exact integer matrix computation in
\texttt{code/D\_exact\_checks.py}; the Smith form and both invariant
spaces are printed in its output. The integral index is proved in Paper~I, Theorem~7.1; the count uses the displayed quotient just computed.
\end{proposition}

The monodromy invariants, coinvariants and resulting corrected component count are contained in \citet[Lemma~2.7, Section~10.5, p.~91, and Appendix~A]{AC2026}. The exact matrix calculation below is retained as a re-derivation in the common marking.



